\chapter{IA integree a l'application}
\label{ch:app}

Ce chapitre decrit \textbf{comment l'intelligence artificielle est utilisee
dans le produit Promptly} (cote application), et non comment le code a ete
genere avec Cursor / Claude / Grok. Deux familles d'outils coexistent :
\begin{itemize}
  \item \textbf{Ollama} : modele de langage local (LLM) ;
  \item \textbf{APIs externes} et \textbf{microservices ML} : services
        specialises (moderation, emotion, matching, prediction).
\end{itemize}

Les cles API sont stockees dans
\texttt{application-local.properties} (fichier gitignore), jamais dans le
depot Git. Documentation detaillee : \texttt{docs/apis-et-tokens.md}.

\section{Vue d'ensemble des fonctionnalites}

\begin{table}[H]
\centering
\begin{tabular}{@{}p{3.8cm}p{4cm}p{5.2cm}@{}}
\toprule
Fonctionnalite & Technologie & Role \\
\midrule
Resume PDF & Ollama + PDFBox & Resume automatique d'un cours / chapitre \\
Flashcards & Ollama + PDFBox & 4 cartes pedagogiques depuis le PDF \\
Generation de quiz & Ollama & 10 QCM a partir du titre / cours / PDF \\
Assistant IA & Ollama & Chatbot general dans l'application \\
Emotion reclamations & Hugging Face & Ton (colere, neutre, \ldots) pour l'admin \\
Censure forum & Sightengine & Insultes remplacees par \texttt{***} \\
Auto-correct forum & TextGears & Orthographe / grammaire a la publication \\
Matching CV / offre & Google Gemini & Score competences + experience \\
GUESS / KNOW quiz & ml-integrity (:8090) & Reponse au hasard vs connue \\
Type / urgence & ml-complaints (:8091) & Labels admin sur reclamations \\
\bottomrule
\end{tabular}
\caption{Fonctionnalites d'IA integrees a Promptly}
\end{table}

\section{Ollama --- LLM local}

\subsection{Pourquoi Ollama}

Ollama permet d'executer un modele de langage \textbf{sur la machine de
developpement} (\texttt{http://localhost:11434}), sans cout API pour les
fonctions pedagogiques centrales. Les donnees des PDF et des conversations
restent locales. Modele utilise : \texttt{llama3.2:3b} (ou equivalent
installe via \texttt{ollama pull}).

Extraction de texte des PDF : bibliotheque Java \textbf{Apache PDFBox}
dans le backend Spring Boot, avant envoi du texte a Ollama.

\subsection{Resume automatique de PDF}

Sur la fiche d'un \textbf{cours} ou d'un \textbf{chapitre} ayant un PDF :
\begin{enumerate}
  \item l'utilisateur clique sur << Resumer automatiquement >> ;
  \item le backend extrait le texte du PDF (tronque si trop long) ;
  \item Ollama produit un resume en francais ;
  \item le resume s'affiche sous la section PDF.
\end{enumerate}
Endpoints typiques :
\texttt{POST /api/ai/courses/\{id\}/summarize-pdf} et equivalent chapitre.
Limitation : un PDF scanne (image seule, sans texte selectable) ne peut pas
etre resume.

\subsection{Flashcards}

Meme flux PDF $\rightarrow$ texte $\rightarrow$ Ollama, avec une consigne
demandant \textbf{4 flashcards} (recto = titre / question, verso = tip).
Bouton << Generer 4 flashcards >> ; l'utilisateur retourne les cartes dans
l'interface Angular. Le backend parse la reponse JSON d'Ollama (avec
fallbacks si le format est imparfait).

\subsection{Generation de quiz (10 questions)}

Dans le formulaire << Nouveau quiz >>, apres titre, description et choix
du cours (chapitre optionnel) :
\begin{enumerate}
  \item clic sur << Generer 10 questions avec IA >> ;
  \item Ollama recoit titre, description et contenu du cours / chapitre / PDF ;
  \item 10 QCM sont proposes ; l'enseignant peut les modifier avant creation.
\end{enumerate}
La generation peut prendre 30 a 90 secondes selon la machine.

\subsection{Assistant IA (chatbot)}

Page \texttt{/ia} : chatbot general de la plateforme.
\begin{itemize}
  \item Backend : \texttt{POST /api/ai/chat} avec historique de conversation ;
  \item Frontend : bulles de chat, suggestions de questions, envoi par Entree ;
  \item Pre-requis : Ollama demarre et modele telecharge.
\end{itemize}
L'assistant repond a des questions sur Java, Spring, Angular, le fonctionnement
de Promptly, etc. Ce n'est pas un remplacement du professeur : c'est une aide
contextuelle integree a l'application.

\section{APIs externes}

\subsection{Hugging Face --- emotion des reclamations}

\textbf{Usage.} Analyser le \textbf{ton} d'une reclamation (En colere, Content,
Triste, Inquiet, Surpris, Neutre) a partir du texte objet + description.

\textbf{Qui voit le resultat.} Uniquement l'\textbf{administrateur} (panneau
Analyse ML). L'etudiant qui depose la reclamation ne voit pas l'emotion.

\textbf{Technique.} Appel Inference API Hugging Face depuis le backend Spring
(\texttt{ComplaintEmotionService}). Token :
\texttt{promptly.emotion.hf-token}. Si l'API est indisponible, un
\textbf{lexique local} FR/EN sert de secours.

\subsection{Sightengine --- moderation du forum}

\textbf{Usage.} Detecter les insultes / gros mots dans les \textbf{posts} et
\textbf{commentaires} du forum, puis les remplacir par \texttt{***} (plutot
que de rejeter le message).

\textbf{Technique.} API Sightengine (plan gratuit) avec \texttt{api-user} et
\texttt{api-secret} dans \texttt{application-local.properties}. Service Java :
\texttt{ProfanityFilterService}. Sans cles : liste locale de mots.

\subsection{TextGears --- auto-correction du forum}

\textbf{Usage.} Corriger l'orthographe et la grammaire \textbf{a la
publication} d'un post ou d'un commentaire (avant la censure Sightengine).

\textbf{Technique.} API TextGears avec cle
\texttt{promptly.spellcheck.api-key} ; langue configurable
(\texttt{fr-FR} ou \texttt{en-US}). Service : \texttt{SpellCorrectService}.
Le professeur exigeait une API avec token : TextGears a ete retenu apres
essais d'autres fournisseurs (LanguageTool Premium, GrammarBot ferme, etc.).

\subsection{Google Gemini --- matching CV / offre}

\textbf{Usage.} Dans le module \textbf{Recrutement}, apres depot d'un CV PDF
sur une offre contenant des \textbf{competences} (et eventuellement une
experience) :
\begin{enumerate}
  \item extraction du texte du CV (PDFBox) ;
  \item Gemini compare le CV aux competences de l'offre ;
  \item le backend stocke un \textbf{score Match \%}, les competences
        trouvees / manquantes, et un indicateur d'experience.
\end{enumerate}
Visible par l'\textbf{admin} sur la fiche de candidature. Cle gratuite AI Studio :
\texttt{promptly.cv-match.api-key}. Sans cle : matching lexical (mots-cles)
en secours. Hugging Face n'est \textbf{plus} utilise pour cette fonction
(uniquement pour l'emotion).

\section{Microservices ML locaux (Python)}

\subsection{ml-integrity (port 8090) --- GUESS / KNOW}

Lorsqu'un etudiant passe un quiz (une question a la fois, avec minuteurs),
le service Python \texttt{ml-integrity} (scikit-learn, RandomForest) predit
si la reponse est :
\begin{itemize}
  \item \texttt{GUESS} : reponse au hasard ;
  \item \texttt{KNOW} : reponse probablement connue.
\end{itemize}
Features utilisees : temps passe, changements de reponse, etc.
\textbf{Seul l'enseignant} voit ces labels (Detail + analyse ML). Une
reponse fausse est forcee en \texttt{GUESS}. Si le service est arrete, Spring
applique une heuristique locale.

\subsection{ml-complaints (port 8091) --- type, urgence, priorite}

A chaque nouvelle reclamation, le service \texttt{ml-complaints} predit :
\begin{itemize}
  \item \textbf{type} (technique, pedagogique, compte, \ldots) ;
  \item \textbf{urgence} (urgent / non urgent) ;
  \item \textbf{priorite} (haute / moyenne / \ldots).
\end{itemize}
L'etudiant ne remplit que l'objet et la description ; les labels sont
reserves a l'admin. La classification s'appuie sur le \textbf{texte} (pas
uniquement sur des mots-cles fixes). Emotion = service HF separe (voir
ci-dessus), pas ce microservice.

\section{Synthese}

\begin{itemize}
  \item \textbf{Ollama} couvre le coeur pedagogique : resume, flashcards,
        generation de quiz, assistant chatbot --- gratuit et local.
  \item \textbf{APIs externes} couvrent des besoins specialises : emotion
        (HF), moderation (Sightengine), orthographe (TextGears), matching CV
        (Gemini).
  \item \textbf{ML Python} couvre des predictions metier : integrite des
        reponses au quiz, tri des reclamations.
  \item En cas d'indisponibilite d'une API ou d'un service ML, des
        \textbf{fallbacks} locaux limitent la panne de l'application.
\end{itemize}
